% §16.1 What This Chapter Is For
%
% Monograph register: state the three failure modes and the chapter's
% honesty posture. Reframed away from the per-exercise deferral
% accounting of the textbook draft (those exercises are removed in the
% monograph, so the old "chapter 1 exercise 10 asked..." references would
% dangle); the phenomena are stated directly instead. Cross-refs
% ax:element-indep, ax:block-rates-indep, sec:axiom-1, sec:axiom-2,
% sec:varying-rates, sec:correlated-noise, sec:adversarial-inputs,
% sec:where-framework-helps. DEFINED: sec:what-this-chapter-is-for.

\section{What This Chapter Is For}
\label{sec:what-this-chapter-is-for}

The preceding chapters wrote the affirmative theory of the Bernoulli
framework and, at several points, set aside a phenomenon that does not
fit the two axioms cleanly: a rate that drifts, a pair of queries that
share a noise source, an input chosen by an adversary. Qualifying each
result in place would have buried the positive content under footnotes;
the negative space is collected here instead. This chapter maps it,
naming for each failure the axiom it violates and the analytical tool
that replaces the broken machinery, and \Cref{sec:where-framework-helps}
recovers the part of the framework that survives.

\subsection*{Three failure modes}

The phenomena set aside earlier fall into three modes, each taken up in
its own section.

\textbf{Varying rates over time (\Cref{sec:varying-rates}).} The error
rate $\fprate$ (or $\fnrate$) is not a stable property of the channel. It
drifts with the query index, because the structure changes (a Bloom
filter filled past capacity), the environment changes (a sensor's noise
grows with temperature), or the calibration data did not represent the
deployment data (classifier drift). This violates
\Cref{ax:block-rates-indep} (Axiom 2): the per-block rates are no longer
stable across queries. The channel matrix $Q$ becomes a sequence $(Q_t)$,
and the iid concentration tools of chapter 11 give way to their
martingale counterparts.

\textbf{Correlated noise across queries (\Cref{sec:correlated-noise}).}
Two queries to the same structure share error sources: the same query
repeated is deterministic, two queries to overlapping Bloom-filter bits
are positively correlated, sequential draws from an online learner carry
a temporal dependence. This violates \Cref{ax:element-indep} (Axiom 1):
the error events at distinct inputs are no longer mutually independent.
The factored joint distribution behind every multi-query formula becomes
an inequality, and the composition theorem of chapter 3 picks up
correction terms whose sign tracks the sign of the correlation.

\textbf{Adversarial inputs (\Cref{sec:adversarial-inputs}).} The query
stream is not iid; the queries are chosen by an adversary, or an adaptive
learner, to expose the worst case. Both axioms can break: the adversary's
$t$-th query is conditional on the responses to the earlier ones (so the
inputs are not independent), and the worst-case query sits at the tail of
whatever per-block rate distribution the channel admits (so the single
$\fprate$ the analyst computed no longer describes the queries that
actually arrive). The right tools are worst-case bounds, cryptographic
hash families, and trapdoor primitives, not concentration inequalities.

\subsection*{Posture}

The chapter is written under a deliberate honesty discipline. Where the
framework genuinely does not apply, it says so, names the violated axiom,
and points to the replacement tool; it does not stretch the axioms to
cover cases where they manifestly fail. \Cref{sec:where-framework-helps}
pushes back against pure relinquishment by cataloguing what survives: the
perfect-filter line of chapter 6 still classifies systems by their
rate-plane position when the rates are time-varying, the channel-matrix
view of chapter 9 still applies query by query when the matrix is not
stable, and the PPV/NPV machinery of chapter 7 still holds at any single
query even when the prevalence drifts. These are partial successes,
catalogued honestly rather than oversold.
